% Źródła: Eves s. 347--353 (§8.5), 353 nn. (§8.6); Wikipedia: Poincaré disk model, Henri Poincaré, Felix Klein.
% https://en.wikipedia.org/wiki/Poincar%C3%A9_disk_model
% Uwaga: Wikipedia przypisuje metrykę dysku już Beltramiemu (1868); nazwa pochodzi od Poincarego (ok. 1882).

\section{Poincaré i niesprzeczność} % lang-pl
\section{Poincaré e la coerenza} % lang-it
\label{sekcja_poincare}

Tego, co Beltrami zrobi po cichu, Poincaré dokona głośno i przy okazji zupełnie innej sprawy. % lang-pl
W tej sekcji zobaczymy, skąd Poincaré weźmie swój model, jak on wygląda i co dokładnie z niego wynika: niesprzeczność geometrii Łobaczewskiego względem euklidesowej oraz niezależność piątego postulatu. % lang-pl
Ciò che Beltrami avrà fatto in silenzio, Poincaré lo farà a voce alta e in occasione di tutt'altra questione. % lang-it
In questa sezione vedremo da dove Poincaré prenderà il suo modello, come si presenta e che cosa ne discende esattamente: la coerenza della geometria di Lobachevskij rispetto a quella euclidea e l'indipendenza del quinto postulato. % lang-it

Henri Poincaré urodzi się 29 kwietnia 1854 roku w Nancy, a umrze 17 lipca 1912 roku w Paryżu. % lang-pl
Model dysku wykorzysta w 1882 roku w pracach o funkcjach automorficznych (fuchsowskich), w których będzie nieformalnym rywalem Kleina od 1881 roku; szersze uznanie zdobędzie on dopiero po wydaniu w 1905 roku jego filozoficznej książki \emph{La science et l'hypothèse}. % lang-pl
Henri Poincaré nascerà il 29 aprile 1854 a Nancy e morirà il 17 luglio 1912 a Parigi. % lang-it
Userà il modello del disco nel 1882 nei lavori sulle funzioni automorfe (fuchsiane), in cui sarà dal 1881 un amichevole rivale di Klein; un più ampio riconoscimento lo otterrà solo dopo la pubblicazione, nel 1905, del suo libro filosofico \emph{La science et l'hypothèse}. % lang-it
\index[persons]{Poincaré, Henri}%

\begin{definition}
[model Poincarego] % lang-pl
[modello di Poincaré] % lang-it
	Ustalmy okrąg $\Sigma$ (\emph{okrąg podstawowy}). % lang-pl
	Punktami są punkty wnętrza $\Sigma$; prostymi -- części wnętrza $\Sigma$ dowolnych „okręgów'' (okręgów lub prostych) prostopadłych do $\Sigma$; długość odcinka $AB$ to $\log(AB, TS)$, gdzie $S$, $T$ są przecięciami „okręgu'' zawierającego $AB$ z $\Sigma$; miara kąta to miara euklidesowa kąta między „okręgami''. % lang-pl
	Fissiamo una circonferenza $\Sigma$ (\emph{circonferenza fondamentale}). % lang-it
	I punti sono i punti dell'interno di $\Sigma$; le rette sono le parti interne a $\Sigma$ di «circonferenze» qualsiasi (circonferenze o rette) ortogonali a $\Sigma$; la lunghezza del segmento $AB$ è $\log(AB, TS)$, dove $S$, $T$ sono le intersezioni con $\Sigma$ della «circonferenza» che contiene $AB$; la misura di un angolo è la misura euclidea dell'angolo tra le «circonferenze». % lang-it
\index{model!Poincarego} % lang-pl
\index{modello!di Poincaré} % lang-it
\end{definition}

Eves \cite[s. 347--348]{eves1_1972}. % lang-pl
Eves \cite[p. 347--348]{eves1_1972}. % lang-it

Dysk ten jest \emph{konforemny} (zachowuje kąty), a ten sam dysk pojawił się już w pracy Beltramiego z 1868 roku. % lang-pl
Kilkadziesiąt lat później M.~C. Escher użyje go w serii \emph{Circle Limit} (1958--1960) po rozmowach z Harold Coxeterem. % lang-pl
Il disco è \emph{conforme} (conserva gli angoli), e lo stesso disco compariva già nel lavoro di Beltrami del 1868. % lang-it
Qualche decennio dopo M.~C. Escher lo userà nella serie \emph{Circle Limit} (1958--1960) dopo le conversazioni con Harold Coxeter. % lang-it
\index[persons]{Escher, Maurits Cornelis}%
\index[persons]{Coxeter, Harold}%

\begin{proposition}
	Jeśli geometria euklidesowa (na płaszczyźnie) jest niesprzeczna, to geometria Łobaczewskiego (na płaszczyźnie) też jest niesprzeczna. % lang-pl
	Se la geometria euclidea del piano è coerente, allora anche la geometria di Lobachevskij del piano è coerente. % lang-it
\end{proposition}

Eves \cite[s. 347-353]{eves1_1972}. % lang-pl
Eves \cite[p. 347-353]{eves1_1972}. % lang-it

Eves bierze za punkt wyjścia układ aksjomatów Hilberta, w którym zamienia aksjomat równoległych (IV-1) na aksjomat Łobaczewskiego IV$'$-1: przez punkt spoza prostej przechodzą co najmniej dwie proste, które jej nie przecinają \cite[s. 347]{eves1_1972}. % lang-pl
Sprawdza wszystkie aksjomaty grup I--III (łączności, porządku, przystawania), IV$'$-1 i V$'$-1 (ciągłość) w modelu: większość jest oczywista, a przystawanie wymaga lematu o niezmienniczości długości względem inwersji w okręgu prostopadłym do $\Sigma$ \cite[s. 348--352]{eves1_1972}. % lang-pl
Eves parte dal sistema di assiomi di Hilbert, in cui sostituisce l'assioma delle parallele (IV-1) con l'assioma di Lobachevskij IV$'$-1: per un punto esterno a una retta passano almeno due rette che non la intersecano \cite[p. 347]{eves1_1972}. % lang-it
Verifica nel modello tutti gli assiomi dei gruppi I--III (connessione, ordine, congruenza), IV$'$-1 e V$'$-1 (continuità): la maggior parte è ovvia, e la congruenza richiede un lemma sull'invarianza della lunghezza rispetto all'inversione in una circonferenza ortogonale a $\Sigma$ \cite[p. 348--352]{eves1_1972}. % lang-it

Dowód ma jeszcze drugą stronę: skoro w modelu (będącym częścią geometrii euklidesowej) spełnione są wszystkie postulaty euklidesowe oprócz piątego, a piąty zawodzi, to piąty postulat \emph{nie} wynika z pozostałych -- jest od nich niezależny \cite[s. 352]{eves1_1972}. % lang-pl
Eves zaznacza też, że sama niesprzeczność geometrii euklidesowej sprowadza się ostatecznie do niesprzeczności liczb rzeczywistych, która pozostaje otwartym pytaniem \cite[s. 352--353]{eves1_1972}. % lang-pl
La dimostrazione ha anche un altro lato: poiché nel modello (che fa parte della geometria euclidea) valgono tutti i postulati euclidei tranne il quinto, e il quinto cade, il quinto postulato \emph{non} discende dagli altri -- è indipendente da essi \cite[p. 352]{eves1_1972}. % lang-it
Eves nota anche che la stessa coerenza della geometria euclidea si riduce in ultima analisi alla coerenza dei numeri reali, che rimane una questione aperta \cite[p. 352--353]{eves1_1972}. % lang-it

Model służy też do obliczeń: to, co trudno wyprowadzić w geometrii Łobaczewskiego, często łatwo wyprowadzić z jego odpowiednika w geometrii euklidesowej, a nawet cała trygonometria hiperboliczna wynika z modelu \cite[s. 353 nn.]{eves1_1972}; Eves opublikuje ją razem z V.~E. Hoggattem w \emph{American Mathematical Monthly} w 1951 roku. % lang-pl
Il modello serve anche per i calcoli: ciò che è difficile ricavare nella geometria di Lobachevskij si ricava spesso facilmente dal suo corrispondente nella geometria euclidea, e persino tutta la trigonometria iperbolica discende dal modello \cite[p. 353 e segg.]{eves1_1972}; Eves la pubblicherà con V.~E. Hoggatt nell'\emph{American Mathematical Monthly} nel 1951. % lang-it
\index[persons]{Hoggatt, Verner Emil}%
\index[persons]{Eves, Howard}%

\todofoot{Zakres §8.6 Evesa (deduktywne twierdzenia z modelu Poincarego) przeczytałem tylko do s. 354; strony w tym akapicie wymagają uzupełnienia po pełnej lekturze.} % lang-pl
\todofoot{Ho letto il §8.6 di Eves (deduzioni dal modello di Poincaré) solo fino a p. 354; le pagine di questo paragrafo vanno completate dopo una lettura integrale.} % lang-it
